% Two real per-patient Suppression Ratio (SR) signals -- a low-score
% survivor and a high-score non-survivor -- with the two filter ranges
% used in this paper drawn as shaded bands:
%   [75, 100) (red)   -- top range, Non-survivor signature
%                        (cell [-1, 0] of GeneralRangeFilter(4, 1))
%   [ 0,  25) (green) -- bottom range, Survivor signature
%                        (cell [ 0, 0] of GeneralRangeFilter(4, 1))
% Both projections are forwarded to single-threshold linear classifiers
% in Section~\ref{sec:results}. The per-patient filter score reported
% in each panel title is the fraction of the recording inside the
% corresponding band.
%
% Data:
%   data/sr_example_nonsurvivor.csv  -- PUH-2012-213, top-range = 0.702
%   data/sr_example_survivor.csv     -- PUH-2015-133, top-range = 0.162
% produced by filter-repository-python/export_sr_example_signals.py.
\usepgfplotslibrary{groupplots}
\begin{tikzpicture}
    \begin{groupplot}[
        group style={
            group size=1 by 2,
            vertical sep=1.6cm,
            xlabels at=edge bottom,
        },
        width=\linewidth,
        height=3.8cm,
        xmin=0, xmax=1,
        ymin=0, ymax=100,
        ytick={0, 25, 50, 75, 100},
        extra y ticks={25, 75},
        extra y tick labels={},
        extra y tick style={grid=major, grid style={dashed, gray!60, thick}},
        enlarge x limits=false,
        xlabel={Normalised recording time},
        ylabel={Suppression Ratio},
        grid=major,
        grid style={dashed, gray!25},
        title style={font=\small\bfseries, align=center},
        legend columns=2,
        legend style={
            font=\footnotesize, draw=none, fill=none,
            /tikz/every even column/.append style={column sep=0.8em},
        },
    ]

    % ---------- Panel 1: Non-survivor ----------
    \nextgroupplot[title={Non-survivor (PUH-2012-213)}]
        % Top range: Non-survivor signature.
        \fill[red!12] (axis cs:0,75) rectangle (axis cs:1,100);
        % Bottom range: Survivor signature.
        \fill[green!12] (axis cs:0,0)  rectangle (axis cs:1,25);
        \addplot+[mark=none, forget plot, line width=0.4pt, draw=red!50!black]
            table [col sep=comma, x=t, y=sr]
            {data/sr_example_nonsurvivor.csv};
        % All legend entries collected on the second panel below so the
        % shared `legend to name=srlegend' is only emitted once.

    % ---------- Panel 2: Survivor (also collects the shared legend) ----------
    \nextgroupplot[title={Survivor (PUH-2015-133)},
                   legend to name=srlegend]
        \fill[red!12]   (axis cs:0,75) rectangle (axis cs:1,100);
        \fill[green!12] (axis cs:0,0)  rectangle (axis cs:1,25);
        \addplot+[mark=none, forget plot, line width=0.4pt, draw=green!40!black]
            table [col sep=comma, x=t, y=sr]
            {data/sr_example_survivor.csv};
        \addlegendimage{mark=none, thick, draw=red!50!black}
        \addlegendentry{SR signal (Non-survivor)}
        \addlegendimage{mark=none, thick, draw=green!40!black}
        \addlegendentry{SR signal (Survivor)}
        \addlegendimage{area legend, fill=red!12, draw=red!12}
        \addlegendentry{Top range $[75, 100)$}
        \addlegendimage{area legend, fill=green!12, draw=green!12}
        \addlegendentry{Bottom range $[0, 25)$}

    \end{groupplot}
    \node[anchor=north] at ($(group c1r2.south) + (0, -1.0cm)$) {\ref{srlegend}};
\end{tikzpicture}
